\documentclass{article}
\usepackage{tikz}
\usetikzlibrary{angles, quotes}
\usepackage{amssymb}
\usepackage{pgfplots}
\pgfplotsset{width=10cm,compat=1.9}
\usepgfplotslibrary{statistics}
\usepackage[left=2.54cm,right=2.54cm,top=2.54cm,bottom=2.54cm]{geometry}
\usepackage{tasks} %Paquete necesario para la producción de items horizontales para algunos ejercicios de tarea empleando el comando /begin{multicols}}{#}.../end{multicols} para ayudarnos a reducir las páginas
\usepackage{pdflscape} %comando que permite cambiar a orientación horizontal las hojas%
\usepackage{soul} %Paquete para permitir la compilación de acentos usando el comando {\hl{}}}\\
\sethlcolor{yellow(munsell)} %comando para definir el color para el subrayado usando el comando \hl
\usepackage[utf8]{inputenc}
\usepackage[spanish]{babel}
\usepackage{icomma}
\usepackage{ragged2e}
\usepackage{siunitx}
\usepackage{url}
\usepackage[colorlinks=true, urlcolor=blue,  linkcolor=black, citecolor=green]{hyperref}
\usepackage{pdfpages}
\usepackage{blindtext} %Paquete que genera texto ficticio (dummy text) con el comando: \blindtext
\usepackage{cancel} %in the preamble gives you four different modes of striking through: \cancel{text to cancel} draws a diagonal line (slash) through its argument, \bcancel{text to cancel} uses the negative slope (a backslash), \xcancel{text to cancel} draws an X (actually \cancel plus \bcancel), \cancelto{〈value〉}{〈expression〉} draws a diagonal arrow through the 〈expression〉pointing to the 〈value〉 (math-mode only)
\usepackage{csquotes}
\usepackage{afterpage}
\usepackage{parskip} 
\usepackage{float}
\usepackage{enumitem}
\usepackage{multicol}%Paquete que permite la creación aislada de columnas de texto. De esta forma se reduce la cantidad de páginas en nuestro documento
\usepackage{lipsum} %Paquete que genera texto de relleno al igual que \usepackage{blindtext}, se usa con el comando \lipsum[#-#] los #(hashtags) delimitan el número de páginas que deseamos ocupar del paquete lipsum para usarlos como relleno. 

% Definir el comando personalizado para la nota
\newcommand{\mynote}[1]{%
\begin{tikzpicture}[baseline=-0.75ex]
\node[draw, circle, fill=Apple Green!60, inner sep=2pt] (note) {\textbf{Nota:}};
\end{tikzpicture}%
\ \textit{#1}%
}

\newenvironment{Figura}
{\par\medskip\noindent\minipage{\linewidth}}
{\endminipage\par\medskip}
\usepackage{caption}
\usepackage[
backend=biber,
sorting=none,
url=true
]{biblatex} %bibliografía
\addbibresource{biblio.bib}
% Establecer el espacio entre las entradas de la bibliografía
\setlength{\bibitemsep}{1\baselineskip} % Puedes ajustar el valor según tus preferencias
\usepackage{amsmath,amsthm,amssymb,amsfonts}
\usepackage{pifont} %Permite la compilación del comando \xmark: es el símbolo de la tachita
\usepackage{mathtools} %permite la compilación de símbolos para matrices
\usepackage{empheq} %Paquete que se relaciona con \usepackage[most]{tcolorbox} para la creación de cuadros/cajas de colores para delimitar los resultados matemáticos o líneas de texto usando la línea de comando: \begin{empheq}[box={\mymath[colback=orange(webcolor)!70,drop lifted shadow]}]{equation*} \end{empheq}
\usepackage[most]{tcolorbox} %Paquete que se relaciona con \usepackage{empheq} para la creación de cuadros/cajas de colores para delimitar los resultados matemáticos o líneas de texto
\renewcommand{\qedsymbol}{$\blacksquare$}
\usetikzlibrary{positioning,decorations.pathreplacing} %Paquete que permite la compilación de llaves para cuadros sinópticos (brace diagram) 
\usepackage{schemata} %The schemata package is designed just for these brace diagrams
\usepackage{booktabs}

\newcommand\AB[2]{\schema{\schemabox{#1}}{\schemabox{#2}}} %comando o paquete necesario para crear cuadros sinópticos

\newtcbox{\mymath}[1][]{%
nobeforeafter, math upper, tcbox raise base,
enhanced, colframe=blue!30!black,
colback=blue!30, boxrule=1pt,
#1} %Comando importante para encasillar los resultados matemáticos en cajas de diferentes colores, se relaciona con los paquetes \usepackage{empheq} y \usepackage[most]{tcolorbox} 

\newenvironment{sysmatrix}[1]
{\left(\begin{array}{@{}#1@{}}}
{\end{array}\right)}
\newcommand{\ro}[1]{%
\xrightarrow{\mathmakebox[\rowidth]{#1}}%
}
\newlength{\rowidth}% row operation width
\AtBeginDocument{\setlength{\rowidth}{3em}}  %Comando importante para laproducción de lineas de operaciones entre matrices, método de Gauss o Gauss Jordan se relaciona con \newenvironment{sysmatrix}

%formato para cambiar el horario a español
\usepackage[spanish]{datetime2}
\DTMsetdatestyle{spanish}

\renewcommand{\today}{\DTMdisplaydate{\the\year}{\the\month}{\the\day }{-1}}
%%%%%%%%%%%%%%%%%%%%%%%%%%%%%%%%%%%%%%%%%

\usepackage{datetime}
\newcommand{\mycurrenttime}{\xxivtime}

%%%%%%%%%%%%%%%%%%%%%%%%%%%%%%%%%%%%%%%%%%%%%%%%%%%%%%%%%%%%%%%%%%%%%%%%%%%%%%%%%%%%%%%%%%%%%%%%%%%%%%%%%%%%%%%%%%%%%%%%%%%%%%%%%%%%%%%%%%%%%%%%%


%%%%%%%%%%%%%Caja de color para el título%%%%%%%%%%%%%%%%%%%%%%%%%%%%%%%%%%%%%%%%%%%%%%%%%%%%%%

\definecolor{myframecolor}{RGB}{1, 45, 75} % Prussian Blue
\definecolor{myboxcolor}{RGB}{0, 107, 92} % Apple Green

\newtcolorbox{mybox}{
enhanced,
colback=myboxcolor!7, % Color de fondo del cuadro
colframe=myframecolor, % Color del marco del cuadro
arc=0pt,
boxrule=1pt,
borderline west={2mm}{-10mm}{myframecolor}, % Borde en el lado izquierdo
sharp corners=southwest,
width=\linewidth,
before=\par\vspace{\bigskipamount}, % Espacio antes del cuadro
after=\par\vspace{\bigskipamount} % Espacio después del cuadro
}


\newcommand{\euler}{e} %Comando para producir letra e de euler
\newenvironment{remark}{\par\vfill\footnotesize % Vertical white space above the remark and smaller font size
\begin{list}{}{
	\leftmargin=80pt % Indentation on the left
	\rightmargin=60pt}\item\ignorespaces % Indentation on the right
\makebox[-2.5pt]{\begin{tikzpicture}[overlay]
		\node[draw=Horizon!60,line width=2.5pt,circle,fill=Horizon!25,font=\sffamily\bfseries,inner sep=4pt,outer sep=0pt] at (-19pt,5pt){\textcolor{Horizon}{Nota}};\end{tikzpicture}} % Blue Nota in a circle
\advance\baselineskip -1pt}{\end{list}\vskip5pt} % Tighter line spacing and white space after remark
\usepackage{graphicx}

\usepackage{titling}

\input{colores.tex}

%%Producir signos de cita textual``Asignatura''

\newcommand{\xmark}{\ding{55}} %%%comando que genera la tachita

\newenvironment{MyColorPar}[1]{%
\leavevmode\color{#1}\ignorespaces%
}{%
}%

%%%%%%%%%%%%%%%%%%% Título de la tarea, Nombre de alumno

\title{Clase \# 46 - Mecánica Teórica}
\author{\emph{JA} $\boldsymbol{\mid}$}
\date{\today}

%%%%%%%%%%%%%%%%%%% Datos de la Materia

\usepackage{fancyhdr}
\fancypagestyle{plain}{%  the preset of fancyhdr 
\fancyhf{} % clear all header and footer fields
\fancyfoot[R]{
\includegraphics[width=2cm]{LogoFCFMBUAP (1).png}}
\fancyfoot[L]{{\bfseries{\thedate{}}} a las {\bfseries{\mycurrenttime{}}} horas{} (GMT-6, H. Puebla de Zaragoza, Pue)}
\fancyhead[L]{Mecánica teórica}
\fancyhead[R]{\theauthor}
}
\makeatletter
\def\@maketitle{%
\newpage
\null
\vskip 1em%
\begin{center}%
\let \footnote \thanks
{\LARGE \@title \par}%
\vskip 1em%
%{\large \@date}%
\end{center}%
\par
\vskip 1em}
\makeatother

\usepackage{lipsum}  
\usepackage{cmbright}


%%%%%%%%%%%%%%%%%%%%%%%%%%%%%%%%%%%%%%%%%%%%%%%%%%%%%%%%%%%%%%%%%%%%%%%%%%%%%%%%%%%%%%%%%%%%%%%%%%%%%%%%%%%%%%%%%%%%%%%%%%%%%%%%%%%%%%%%%%%%%%%%%%%%%%%%%%%%%%%%%%%%%%%%%%%%%%%%%%%%%%%%%%%%%%%%%%%%%%

\begin{document}



\maketitle


%%%%%%%%%%%%%%%%%%% Datos del profesor, campus y aula

\begin{mybox}
\noindent\begin{tabular}{@{}ll}
	{\bfseries{Profesor}} & Dr. Gilberto Silva Ortigoza\\
	\textcolor{prussianblue}{\bfseries{Campus}} / \textcolor{trueblue}{\bfseries{Facultad}}  & \textcolor{prussianblue}{\bfseries{C.U. BUAP}} /  \textcolor{trueblue}{\bfseries{FCFM}} \\
	\textcolor{prussianblue}{\bfseries{Edificio}} /  \textcolor{trueblue}{\bfseries{EMA4}}    & 	\textcolor{prussianblue}{\bfseries{FM9}} /  \textcolor{trueblue}{\bfseries{302}}   
\end{tabular} 
\end{mybox} \vspace{1cm}

\tableofcontents

\vspace{2cm}

\section{Clase del día de hoy}

Partimos del hamiltoniano:

\[  H= H(q,p,t) = \frac{p^{2}}{m}  \]


¿Cómo funciona la teoría de Hamilton-Jacobi? 

Introducimos una función generatriz de tipo $F_{2}$

Tenemos:


\subsection{Función generatriz de tipo $F_{2}$}


Tenemos:


\[ F_{2} = F_{2}(q,P) = S(q,P)  \]



Tenemos:


\[ p = \frac{\partial S}{\partial q}; Q=\frac{\partial S}{\partial P}   \]


SI depende explicitamente dle tiempo, nosotros tendríamos:


\[  K = H + \frac{\partial S}{\partial t}    \]


Tenemos:


\[ K(Q,P,t) = H(q,p,t) + \frac{\partial S}{\partial t}(q,P,t) \]



Tenemos:


\[ \dot{Q} = \frac{\partial K}{\partial P} =0 \]


\[  \dot{P} = - \frac{\partial K}{\partial Q} =0  \]


Tenemos también:

\[ Q(t) = Q_{0} \]

\[ P(t) = P_{0} \]

SI nosotros tenemos una función $S$, nosotros podemos llevar a cabo este procedimiento. 

Existe este tipo de función generatriz, esta $S$?

Si la $S$ existe se debe de cumplir esta condición: 


\[ H(q,p,t) + \frac{\partial S}{\partial t} (q,P,t) =0 \]


EL nuevo gamiltoniano sea identicamente $0$. 

Tenemos:


\[ H\left( q, \frac{\partial S}{\partial q}, t  \right) + \frac{\partial S}{\partial t} =0 \]


Es una ecuación diferencial en derivadas parciales y esta recibe el nombre de \textbf{ecuación de Hamilton-Jacobi}.


Tendriamos que resolver la ecuación anterior para ver si la $S$ existe. 


Para poder determinar la transformación canónica, necesitamos conocer a $S$.


Para cada valor de $P$, esta es es solución de la ecuación de \textbf{Hamilton-Jacobi}.

Tenemos una familia de soluciones que viene etiquetada por cada una de los parámetros d ela $P$ mayúscula. Necesitamos una función generatriz que dependa de $q$, $P$ y también el tiempo $t$.


Tendríamos nosotros:


\[ p = p(q,P,t) \]


\[\underbrace{Q = Q(q,P,t)}_{q=q(Q,P,t)} \]

sustituimos $q = q(Q,P,t)$

Y la sustituimos en $p=p(Q,P,t)$


Tenemos que tener una familia de soluciones en la ecuación de \textbf{Hamilton-Jacobi}.

Con $S$ generamos las soluciones de \textbf{Hamilton}. 

Requerimos una familia de soluciones de \textbf{Hamilton-Jacobi}.


\textcolor{trueblue}{\textbf{Integral completa}}, es decir son las familias de soluciones que resuelven la ecuación de \textbf{Hamilton-Jacobi}. 

\[ \boxed{ \frac{ \partial^{2}S }{ \partial q \partial P } \neq 0 } \]

\subsection{Ejemplo:}


\begin{center}
\includegraphics[width=0.8\linewidth]{1.png} % <-- Sustituye esta ruta
\\[2mm]
\textbf{.}
\end{center}	

Podemos pasar de partícula libre en dos dimensiones al problema de Kepler en 2 dimensiones. 

\subsection{Ejemplo de partícula libre: \textcolor{Cinnabar}{\textbf{Lujo de detalle}}}

Tenemos:


\[ H = H(q,p,t) = \frac{p^{2}}{2m} \]

Tenemos las ecuaciones de \textbf{Hamilton}.


\[ \dot{q} = \frac{\partial H}{\partial p} = \frac{p}{m} = \frac{p_{0}}{m} \quad \to \quad q(t) = q_{0} + \frac{p_{0}}{m}t  \]

\[ \dot{p} = - \frac{\partial H}{\partial q} = 0. \quad \to \quad p(t)=p_{0}  \]


Tenemos:

\begin{center}
	\includegraphics[width=0.8\linewidth]{2.png} % <-- Sustituye esta ruta
\\[2mm]
\textbf{.}
\end{center}	


\[ \boxed{ q(t) = q_{0} + \frac{p_{0}}{m}t  } \]


\[  p(t) = p_{0} \]

\[ q(t) =q_{0} \]

\[ \boxed{ q(t) = q_{0} + \frac{p_{0}}{m}t  }  \]


Tendriamos sustituyendo el momento canónico en la \textbf{función de Hamilton Jacobi}:


\[  \boxed{\frac{1}{2m} \left(  \frac{ \partial S  }{ \partial q }  \right)^{2} + \frac{\partial S  }{\partial t} =0 }  \]


POr que al proponer la solución:

\[ S(q,t) = W(q) + T(t) \]

Es una suma, y no un producto.

Tenemos así:


\[ \frac{\partial S}{\partial q} = \frac{dW}{dq}, \quad \frac{\partial S}{\partial t} = \frac{dT}{dt}  \]


Tenemos:


\[ \frac{1}{2m} \left( \frac{dw}{dq}  \right) = - \frac{dT}{dt} = E   \]

\[ \frac{1}{2m} \left(  \frac{dw}{dq} \right)^{2} = E; \quad \frac{dT}{dt} = -E \quad \to \quad T(t)=-Et + C_{0}\]


\[ \frac{dW}{dq}= \pm \sqrt{ 2mE  } \quad \to \quad W = \pm \sqrt{ 2mE  } q +C_{1}  \]


Tenemos:



\[ F_{2} = F_{2}(q,P,t) = S(q,P,t) = \pm \sqrt{2mE  } q - Et +C_{2}   \]

Tenemos:


\[ p = \frac{\partial S}{\partial q} = \pm \sqrt{ 2mE  }  \]


\[ Q = \frac{\partial S}{\partial E} = \pm \frac{ \sqrt{ 2m  }q }{2\sqrt{E}}-t \]


Tenemos:


\[ P = \pm \sqrt{ 2mE } \]

\[ P^{2} = 2mE \]


\[ \boxed{S(q,P,t) = qP - \frac{P^{2}}{2m} t} \]

\[ \frac{ \partial^{2}}{\partial q \partial P} = 1. \quad \textcolor{Cinnabar}{ p = \frac{\partial S}{\partial q} =P; \, Q=\frac{\partial S}{\partial P} = q - \frac{p}{m}t}  \]


\begin{center}
	\includegraphics[width=0.8\linewidth]{3.png} % <-- Sustituye esta ruta
\\[2mm]
\textbf{.}
\end{center}	

\[ p =P \]


\[  Q = q- \frac{P}{m}t \]


Tenemos:


\[q = Q + \frac{p}{m}t  \]

\[ p = P \]


Tenemos:


\[ \boxed{ Q = q_{0}, \quad P=p_{0}  } \]


En las \textbf{transformaciones de Lorentz} tenemos un sistema de referencia inercial.

¿Cuál es la relación en entre una y otra? 

Tenemos que repasar estas definiciones. 

Las transformaciones de Lorentz en realidad son transformaciones muy particulares, las tranformaciones canónicas son mucho más generales.


\subsection{Oscilador armónico}.

Tenemos nosotros:


\[ H(x,p,t) = \frac{ p^{2} }{m} + \frac{1}{2} m\omega^{2}x^{2} \]

Tenemos:


\[ \dot{x} = \frac{\partial H}{\partial p} = \frac{p}{m}  \]


Tenemos:


\[ \dot{p} = - \frac{\partial H}{\partial x} = -m\omega^{2}{x} \]


\[ \dot{x} = \frac{p}{m}\]


\[ \dot{p }= -m\omega^{2}{x} \]


Tenemos:

\[ x(t) = x_{0}\cos(\omega{t}) + \frac{p_{0}}{m\omega} \sen(\omega{t}) \]

\[ p(t) = p_{0}\cos(\omega{t}) - m\omega x_{0}\sen\omega{t}  \]



Realizamos una tranformación canónica con función generatriz $S$. 


Tenemos:


\[ S(x,t,P): \quad K=0 \]


\[ K = H + \frac{\partial S}{\partial t}  = 0 \]

Tenemos:



\[  \frac{1}{2m} \left(  \frac{ \partial S }{\partial x}  \right)^{2} + \frac{1}{2}m\omega^{2}x^{2} + \frac{\partial S}{\partial t} =0    \]


Tenemos:


\[ S(x,t) = W(x) + T(t)  \]


Tenemos:


\[  \frac{1}{2m} \left(  \frac{ dw }{dx} \right)^{2} + \frac{1}{2} m\omega^{2}x^{2} + \frac{dT}{dt} =0  \]


Tenemos:


\[ \frac{1}{2m} \left(  \frac{ dW  }{dx} \right)^{2} + \frac{1}{2} m\omega^{2}x^{2} = - \frac{dT}{dt} \quad \to \quad \frac{dT}{dt} = -E  \]

Tenemos:


\[ \frac{dw}{dx} = \pm \sqrt{ 2mE - m^{2}\omega^{2}x^{2}  } \]


\[ w(x) = \pm \int\limits \sqrt{  2mE-m^{2}\omega^{2}x^{2}  } dx + C_{1} \]


Tenemos:


\[ S(x,t,E) = \pm \int\limits \sqrt{  2mE-m^{2}\omega^{2}x^{2}   } dx -Et +C_{3} \]


Tenemos:


\[ S(x,t,P) = \pm \int\limits \sqrt{ P^{2} - m^{2}\omega^{2}x^{2}  } dx - \frac{P^{2}}{2m}t +C_{3} \]

Con:


\[ T(t) = -Et +C_{0} \]

\[   C_{3} = C_{0} +C_{1}  \]


También:


\[  \frac{1}{2m} \left(   \frac{dw}{dx}  \right)^{2} + \frac{1}{2} m\omega^{2}x^{2} =E    \]



\[ \left(  \frac{dw}{dx} \right)^{2} + m^{2}\omega^{2}x^{2} =2mE=P^{2}  \]

\[ \frac{dw}{dx} = \pm \sqrt{ P^{2} -m^{2}\omega^{2}x^{2}  } \]

Tenemos:


\[ W(x) = \pm \int\limits \sqrt{ P^{2}-m^{2}\omega^{2}x^{2}  }dx \]

\[ E = \frac{P^{2}}{2m} \]

Tenemos:


\[ H = \frac{p^{2}}{2m} + \frac{1}{2}m\omega^{2}x^{2}  \]


\[  p = \frac{\partial S}{\partial x} = \pm \sqrt{ P^{2}-m^{2}\omega^{2}x^{2}  } \]


\[ p^{2} = P^{2}-m^{2}\omega^{2}x^{2}\geq 0 \]

\[p^{2} - m^{2}\omega^{2}x^{2}=P^{2} \]

Tenemos:


\[ m^{2}\dot{x}^{2} = P^{2} -m^{2}\omega^{2}x^{2}\geq 0  \]


\[ P^{2} = m^{2}\omega^{2}x_{0}^{2}  \]

\[ X_{0} = \pm \sqrt{ \frac{P^{2}}{m^{3}\omega^{2}}  } = \pm \frac{P}{m\omega} \]




\end{document}